\documentclass[12pt,a4paper]{article}
\usepackage[utf8]{inputenc}
\usepackage[T1]{fontenc}
\usepackage[brazil]{babel}
\usepackage{amsmath,amssymb}
\usepackage{geometry}
\geometry{margin=2.5cm}
\title{Material de Estudo: Discuss\~ao de Exerc\'icios II}
\author{Princ\'ipio de Modelagem Matem\'atica}
\date{Unidade 18}
\begin{document}
\maketitle

\section*{Objetivo}
Integrar os t\'opicos estoc\'asticos (Unidades 16--17): equa\c{c}\~oes de Langevin, Fokker-Planck e cadeias de Markov em problemas realistas.

\section*{Problemas Integrados}

\subsection*{E1 --- Ornstein-Uhlenbeck: precifica\c{c}\~ao de commodities}
O pre\c{c}o de uma commodity reverte \`a m\'edia: $dP = \theta(\mu - P)\,dt + \sigma\,dW$.
\begin{enumerate}
  \item[$a.$] Interprete cada par\^ametro ($\theta$, $\mu$, $\sigma$).
  \item[$b.$] Escreva a equa\c{c}\~ao de Fokker-Planck para $p(P,t)$.
  \item[$c.$] Qual a distribui\c{c}\~ao estacion\'aria? (Dica: Gaussiana com m\'edia $\mu$ e vari\^ancia $\sigma^2/(2\theta)$.)
  \item[$d.$] Para $\mu=100$, $\theta=0{,}5$, $\sigma=10$: calcule a vari\^ancia estacion\'aria e o intervalo de 95\%.
\end{enumerate}

\subsection*{E2 --- SIR estoc\'astico como cadeia de Markov}
Numa popula\c{c}\~ao de $N=10$ indiv\'iduos, o estado \'e $(S, I)$ com $R = N-S-I$.
Transi\c{c}\~oes poss\'iveis em $\Delta t$:
\begin{itemize}
  \item Infec\c{c}\~ao: $(S, I) \to (S-1, I+1)$ com taxa $\beta SI/N$.
  \item Recupera\c{c}\~ao: $(S, I) \to (S, I-1)$ com taxa $\gamma I$.
\end{itemize}
\begin{enumerate}
  \item[$a.$] Por que $I=0$ \'e um estado absorvente?
  \item[$b.$] Para $\beta=0{,}3$ e $\gamma=0{,}1$, calcule $\mathcal{R}_0 = \beta/\gamma$. A epidemia se propaga?
  \item[$c.$] Liste os estados poss\'iveis se $I(0)=1$, $S(0)=9$.
\end{enumerate}

\subsection*{E3 --- Processo de Poisson para chegadas em banco}
Clientes chegam ao banco como processo de Poisson com taxa $\lambda = 3$ clientes/min.
\begin{enumerate}
  \item[$a.$] Qual a probabilidade de 0 clientes em 1 min? E de 5 clientes?
  \item[$b.$] Qual o n\'umero esperado de clientes em 10 min?
  \item[$c.$] Qual a distribui\c{c}\~ao do tempo at\'e o pr\'oximo cliente? (Exponencial com par\^ametro $\lambda$.)
\end{enumerate}

\section*{Exerc\'icios de Revis\~ao}

\begin{enumerate}
  \item A equa\c{c}\~ao de Fokker-Planck para $dX = -\alpha X\,dt + \sigma\,dW$ \'e:
  \[
  \partial_t p = \alpha\partial_x(xp) + \frac{\sigma^2}{2}\partial_{xx}p.
  \]
  Verifique que $p^*(x) \propto \exp(-\alpha x^2/\sigma^2)$ \'e uma solu\c{c}\~ao estacion\'aria.
  
  \item Para a cadeia de Markov de 3 estados com matriz do exerc\'icio D3 (unidade 17), simule 5 passos partindo do estado 1.
  
  \item Explique com palavras: qual a diferen\c{c}a entre a distribui\c{c}\~ao estacion\'aria de uma cadeia de Markov e a distribui\c{c}\~ao de equil\'ibrio da equa\c{c}\~ao de Fokker-Planck?
  
  \item Um processo de nascimento-morte tem taxa de nascimento $\lambda$ e morte $\mu$. Se $\lambda = \mu$, a popula\c{c}\~ao esperada \'e constante, mas a vari\^ancia cresce com $t$. O que isso implica para previs\~oes de longo prazo?
  
  \item Na deriva gen\'etica ($N=100$, $p_0=0{,}5$): ap\'os $t$ gera\c{c}\~oes, $\text{Var}[p_t] \approx p_0(1-p_0)t/N$. Ap\'os quantas gera\c{c}\~oes a vari\^ancia atinge $0{,}1$?
\end{enumerate}

\section*{Gabarito Parcial}
\textbf{E1d:} $\text{Var}^* = 10^2/(2\times0{,}5) = 100$; IC 95\%: $[100-1{,}96\times10, 100+1{,}96\times10] = [80{,}4, 119{,}6]$.\\
\textbf{E3a:} $P(X=0) = e^{-3} \approx 0{,}050$; $P(X=5) = e^{-3}3^5/120 \approx 0{,}101$.\\
\textbf{Rev.\ 5:} $t = 0{,}1 \times 100 / (0{,}5\times0{,}5) = 40$ gera\c{c}\~oes.

\section*{Refer\^encias}
\begin{itemize}
  \item Material das Unidades 16--17.
  \item Notas de aula da disciplina.
\end{itemize}

\end{document}
